\begin{solapartat}
\punts{0,5} Dibuix: La força elèctrica és repulsiva ja que les càrregues tenen el mateix signe i s'oposa al pes.
\figura[0.12]{sol_fig1.png}

\punts{0,25} L'objecte levita, per tant, està en equilibri, i el sumatori de forces que actuen sobre ell serà zero. Per tant, els mòduls de la força elèctrica i del pes seran iguals: $\vec{F}_e + \vec{P} = 0$ i $|\vec{F}_e| = |\vec{P}|$,

\punts{0,5} Per calcular la càrrega de l'objecte, utilitzem l'expressió de la força elèctrica i la igualem a la del pes:
\[ k\,\frac{q_{obj}\,q_{bar}}{r^2} = mg \]
En què la càrrega de l'objecte i de la barra són iguals: $q_{obj} = q_{bar} = Q$ i, per tant, podem escriure: $k\,\frac{Q\cdot Q}{r^2} = mg$ i, per tant,
\[ Q = \sqrt{\frac{mgr^2}{k}} = \sqrt{\frac{0,01\cdot 9,8\cdot 0,55^2}{8,99\cdot 10^{9}}} = 1,82\cdot 10^{-6}\ C = 1,82\ \mu C\,. \]
Per tant, la càrrega de l'objecte és $q_{obj} = 1,82\ \mu C$.
% DUBTE: la pauta dona la càrrega en valor positiu, tot i que l'enunciat diu que la càrrega de la vareta és negativa (seria -1,82 μC). Es transcriu tal com és.
\end{solapartat}

\begin{solapartat}
\punts{0,15} El camp elèctric en el punt central és la suma dels dos camps.

\punts{0,15} Les càrregues són iguals, $-2\ \mu C$, i les distàncies al punt també, $0,3\ m$, per tant ambdós camps tindran el mateix mòdul. El camp elèctric creat per una càrrega negativa està dirigit cap a la càrrega que el crea, així doncs tindran sentit contrari.
\[ \vec{E} = \vec{E}_{obj} + \vec{E}_{bar} = k\,\frac{q}{r^2}\,\vec{\jmath} - k\,\frac{q}{r^2}\,\vec{\jmath} = 0 \]
\figura[0.13]{sol_fig2.png}
L'alumne ha de calcular o justificar el resultat.

\punts{0,15} El potencial serà la suma dels potencials creats per cada càrrega. Considerant el potencial nul a l'infinit:
\[ V = V_{obj} + V_{bar} = k\,\frac{q}{r} + k\,\frac{q}{r}\,. \]

\punts{0,20} Per tant $V = 8,99\cdot 10^{9}\,\frac{-2\cdot 10^{-6}}{0,3} + 8,99\cdot 10^{9}\,\frac{-2\cdot 10^{-6}}{0,3} = -1,20\cdot 10^{5}\ V$

\textit{Càlcul del treball de la força externa:}

\punts{0,20} Per portar un electró des de l'infinit fins al punt central haurem de fer en tot moment una força contraria al camp elèctric, per tant el treball fet per la força externa és el treball fet per la força elèctrica canviat de signe.

\punts{0,40} El treball realitzat pel camp elèctric és menys l'increment d'energia potencial:
\[ W_E = -\Delta U = -q\cdot\Delta V = 1,602\cdot 10^{-19}\cdot(-1,2\cdot 10^{5} - 0) = -1,92\cdot 10^{-14}\ J \]
Per tant, la força externa ha de fer un treball: $W_{ext} = 1,92\cdot 10^{-14}\ J$

\textit{O, alternativament:}

\punts{0,20} El treball de les forces externes és la variació d'energia potencial: $W_{F\,ext} = \Delta U$

\punts{0,40} Per tant, calculem la variació d'energia potencial al portar la càrrega des de l'infinit al punt central: $\Delta U = q\cdot\Delta V = -1,602\cdot 10^{-19}\cdot(-1,2\cdot 10^{5} - 0) = 1,92\cdot 10^{-14}\ J$

i per tant $W_{ext} = 1,92\cdot 10^{-14}\ J$\,.
\end{solapartat}
